% Part: model-theory
% Chapter: models-of-arithmetic
% Section: non-standard-models

\documentclass[../../../include/open-logic-section]{subfiles}

\begin{document}

\olfileid[pt]{mod}{mar}{nst}
\section{Modelos não padrão}

\begin{explain}
Chamamos uma !!{structure} para $\Lang{L_A}$ de padrão se ela for
isomorfa a~$\Struct{N}$. Se uma !!{structure} não for isomorfa
a~$\Struct{N}$, ela é chamada de não padrão.
\end{explain}

\begin{defn}
Uma !!{structure}~$\Struct{M}$ para $\Lang{L_A}$ é \emph{não padrão}
se não for isomorfa a~$\Struct{N}$. Os !!{element}s $x \in
\Domain{M}$ que são iguais a $\Value{\num{n}}{M}$ para algum $n \in
\Nat$ são chamados de \emph{números padrão} (de $\Struct{M}$); os
demais são chamados de \emph{números não padrão}.
\end{defn}

\begin{explain}
Por \olref[stm]{prop:standard-domain}, toda !!{structure} padrão
para~$\Lang{L_A}$ contém apenas !!{element}s padrão. Consequentemente,
uma !!{structure} não padrão deve conter pelo menos um elemento não
padrão. Na verdade, a existência de um !!{element} não padrão garante
que a !!{structure} seja não padrão.
\end{explain}

\begin{prop}
Se uma !!{structure}~$\Struct{M}$ para $\Lang{L_A}$ contém um número
não padrão, então $\Struct{M}$ é não padrão.
\end{prop}

\begin{proof}
Suponha que não, O.s.\@, suponha que $\Struct{M}$ seja padrão, mas
contenha um número não padrão~$x$. Seja $g\colon \Nat \to \Domain{M}$
um isomorfismo. É fácil ver (por indução em~$n$) que
$g(\Value{\num{n}}{N}) = \Value{\num{n}}{M}$. Em outras palavras, $g$
leva números padrão de~$\Struct{N}$ a números padrão
de~$\Struct{M}$. Se $\Struct{M}$ contém um número não padrão, $g$ não
pode ser !!{surjective}, contrariando a hipótese.
\end{proof}

\begin{prob}
Lembre que $\Th{Q}$ contém os axiomas
\begin{align*}
& \lforall[x][\lforall[y][(\eq[x'][y'] \lif \eq[x][y])]] \tag{$!Q_1$}\\
& \lforall[x][\eq/[\Obj 0][x']] \tag{$!Q_2$}\\
& \lforall[x][(\eq[x][\Obj 0] \lor \lexists[y][\eq[x][y']])] \tag{$!Q_3$}
\end{align*}
Dê !!{structure}s~$\Struct{M_1}$, $\Struct{M_2}$, $\Struct{M_3}$ tais que
\begin{enumerate}
\item $\Sat{M_1}{!Q_1}$, $\Sat{M_1}{!Q_2}$, $\Sat/{M_1}{!Q_3}$;
\item $\Sat{M_2}{!Q_1}$, $\Sat/{M_2}{!Q_2}$, $\Sat{M_2}{!Q_3}$; e
\item $\Sat/{M_3}{!Q_1}$, $\Sat{M_3}{!Q_2}$, $\Sat{M_3}{!Q_3}$;
\end{enumerate}
Evidentemente, basta especificar~$\Assign{\Obj{0}}{M_i}$ e
$\Assign{\prime}{M_i}$ para cada uma.
\end{prob}

\begin{explain}
É bastante fácil especificar !!{structure}s não padrão para
$\Lang{L_A}$. Por exemplo, tome a estrutura com !!{domain}~$\Int$ e
interprete todos os símbolos não lógicos da maneira usual. Como os
números negativos não são valores de $\num{n}$ para nenhum~$n$, essa
estrutura é não padrão. Naturalmente, ela não será um \emph{modelo} da
aritmética no sentido de tornar verdadeiras as mesmas sentenças que
$\Struct{N}$. Por exemplo, $\lforall[x][\eq/[x'][\Obj{0}]]$ é falsa.
Ainda assim, podemos demonstrar com facilidade, usando o teorema da
compacidade, que existem modelos não padrão da aritmética.
\end{explain}

\begin{prop}
Seja $\Th{TA} = \Setabs{!A}{\Sat{N}{!A}}$ a teoria
de~$\Struct{N}$. $\Th{TA}$ tem um modelo não padrão !!{enumerable}.
\end{prop}

\begin{proof}
Expanda $\Lang{L_A}$ com uma nova !!{constant}~$c$ e considere o
conjunto de !!{sentence}s
\[
\Gamma = \Th{TA} \cup \{\eq/[c][\num{0}], \eq/[c][\num{1}],
\eq/[c][\num{2}], \dots\}
\]
Qualquer modelo~$\Struct{M^c}$ de~$\Gamma$ conteria !!a{element}~$x =
\Assign{c}{M}$ não padrão, pois $x \neq \Value{\num{n}}{M}$ para todo
$n \in \Nat$. Além disso, é claro que $\Sat{M^c}{\Th{TA}}$, pois
$\Th{TA} \subseteq \Gamma$. Se transformarmos $\Struct{M^c}$ em uma
!!{structure}~$\Struct{M}$ para $\Lang{L_A}$ simplesmente esquecendo
$c$, seu domínio ainda conterá o elemento não padrão~$x$, e também
$\Sat{M}{\Th{TA}}$. Esta última afirmação é garantida porque $c$ não
ocorre em~$\Th{TA}$. Assim, basta mostrar que $\Gamma$ tem um modelo.

Usamos o teorema da compacidade para mostrar que~$\Gamma$ tem um
modelo. Se todo subconjunto finito de~$\Gamma$ é satisfatível, então
$\Gamma$ também é. Considere um subconjunto finito qualquer
$\Gamma_0 \subseteq \Gamma$. $\Gamma_0$ inclui algumas !!{sentence}s
de~$\Th{TA}$ e algumas da forma~$\eq/[c][\num{n}]$, mas apenas um
número finito delas. Suponha que $k$ seja o maior número tal que
$\eq/[c][\num{k}] \in \Gamma_0$. Defina $\Struct{N_k}$ expandindo
$\Struct{N}$ para incluir a interpretação~$\Assign{c}{N_k} = k+1$.
$\Sat{N_k}{\Gamma_0}$: se $!A \in \Th{TA}$, então $\Sat{N_k}{!A}$,
pois $\Struct{N_k}$ é igual a~$\Struct{N}$ em todos os aspectos,
exceto por~$c$, e $c$ não ocorre em~$!A$. E
$\Sat{N_k}{\eq/[c][\num{n}]}$, pois $n \le k$ e $\Value{c}{N_k} =
k+1$. Logo, todo subconjunto finito de~$\Gamma$ é satisfatível.
\end{proof}

\end{document}
